% !TEX root = journal_0421052099.tex
% =====================================================================
% NOT A DOCUMENT. Packages and styles for journal_0421052099.tex, which
% carries the \documentclass and the \begin{document}; building this
% file directly stops with "Missing \begin{document}".
%
% What elsarticle actually loads is fontenc[T1], graphicx and natbib,
% and nothing else worth relying on. It does NOT load hyperref, amsmath,
% amssymb or url -- \url is undefined until something here provides it,
% which is how the data-availability statement first failed to compile.
% It loads geometry only under the 1p/3p/5p options; under `preprint'
% the measure is the 12pt article default, and geometry is loaded below
% to widen it.
% =====================================================================

% =====================================================================
% FONTS: the same set thesis_book uses, and for the same reasons.
% Changed 2026-09-25 at the user's request ("change the journal to
% exactly match the thesis_book styles"). buetcsepgthesis.sty lines
% 68-86 are the reference: Times for text, Helvetica scaled to 92% for
% sans, Computer Modern left in place for maths, microtype with one
% interface skipped. All four are reproduced below.
%
% This REPLACED \usepackage{lmodern}, which was here for a bug rather
% than for a look, and the bug still has to be answered. elsarticle
% typesets \elsaddress with \footnotesize\itshape inside its
% frontmatter, and under `review' that call reaches font selection with
% an empty size -- "Calculating math sizes for size <>", then a run of
% "Font shape ... in size <> not available" warnings substituting size
% <5>. Bisected to \affiliation: remove it and the warnings vanish, and
% all three affiliation forms (full key-value, organization-only, plain
% text) trigger it equally. What clears them is a font family that is
% scalable, so there is a shape to find at whatever size is asked for;
% lmodern was scalable in both text and maths, which is why it worked.
%
% `times' alone does NOT fix it, and this was measured rather than
% assumed: times sets rmdefault/sfdefault/ttdefault and leaves maths on
% Computer Modern, so the warnings simply move from OT1/cmr to OML/cmm
% and OMS/cmsy -- six of them, plus an underfull box. fix-cm is the
% answer that keeps the thesis's actual font: it declares the CM shapes
% at arbitrary sizes, so maths stays literally Computer Modern, as in
% thesis_book, and there is still a shape at size <>. It must load
% before anything asks for a font, hence first in this file.
%
% Elsevier re-typesets in their own fonts at proof stage, so all of this
% affects the submission and review PDF only.
\usepackage{fix-cm}
\usepackage{times}
% times sets the sans face to Helvetica at full size, but Helvetica's
% x-height is far larger than Times's, so sans text set at the same
% nominal size looks bigger and heavier than the prose around it. This
% paper marks identifiers and system names in sans, and in a paragraph
% carrying several of them the page reads speckled. Scaling Helvetica to
% 92% matches the two faces optically. The conventional pairing
% correction, not a size change: the nominal size is unchanged, so
% nothing about the layout or the marking is affected. Verbatim from
% buetcsepgthesis.sty, so the two documents mark identifiers identically.
\usepackage[scaled=0.92]{helvet}
\usepackage{microtype}
% hyperref (loaded at the foot of this file) redefines \@footnotetext to
% wrap footnotes in a PDF link anchor once hyperfootnotes are active;
% microtype's `footnote' interface patch does a literal-text match
% against the kernel's original \@footnotetext and no longer finds it, so
% it always warns "Unable to apply patch `footnote'" even though no
% footnote has been typeset yet. Skip just that one interface -- all
% other microtype patches (item, toc, eqnum, protrusion and expansion in
% general) are unaffected. Same load order as the thesis: microtype
% before hyperref.
\microtypesetup{nopatch=footnote}

% The measure. Under `preprint' elsarticle leaves the 12pt article
% default, a 131mm by 220mm text block on A4 with 39mm at each side and
% 45mm at the top; with `review' spacing that ran the manuscript to 72
% pages. Elsevier places no layout requirement on an initial submission,
% so the block is widened here to the class's own 3p width (468pt) and
% deepened to 660pt, centred: a 25mm side margin and about 32mm top and
% bottom. The generated figures are \textwidth-relative and follow it;
% their fixed heights do not change.
\usepackage[textwidth=468pt, textheight=660pt, centering]{geometry}

\usepackage{booktabs}
\usepackage{array}
% Algorithm 1 in framework.tex (September 2026). The float package must
% load before hyperref/cleveref so that \Cref{alg:...} resolves as
% "Algorithm"; both are loaded at the end of this file.
\usepackage{algorithm}
\usepackage{algpseudocode}
% Appendix A sets the thirteen prompt templates verbatim, each in a
% lstlisting of this style. basicstyle runs at the start of every
% listing, so \linespread{1}\selectfont there single-spaces the
% listings under `review' the way framework.tex single-spaces
% Algorithm 1: a prompt is code, and 1.5-spaced code reads as prose
% that has lost its paragraphs. The surrounding notes keep the
% document's spacing. breaklines is what keeps the log clean: the
% templates carry lines far longer than the measure.
\usepackage{listings}
\lstdefinestyle{prompt}{
  basicstyle=\linespread{1}\selectfont\footnotesize\ttfamily,
  columns=fullflexible,
  keepspaces=true,
  breaklines=true,
  breakindent=0pt,
  frame=single,
  framerule=0.3pt,
  rulecolor=\color{black!40},
  xleftmargin=4pt,
  xrightmargin=4pt,
}
\usepackage{tikz}
\usepackage{pgfplots}
\pgfplotsset{compat=1.18}
\usetikzlibrary{arrows.meta,positioning,shapes.geometric,calc,fit,
  backgrounds,patterns}

% Semantic figure palette, identical to buetcsepgthesis.sty and to the
% deck: blue for a step that does work, green for supported, vermillion
% for unsupported. The generated figures repeat their own \definecolor
% so they stay self-contained; these are for the hand-drawn ones.
\definecolor{agrNode}{HTML}{0072B2}
\definecolor{agrSup}{HTML}{009E73}
\definecolor{agrUns}{HTML}{D55E00}

% The plot style the generated figures expect. Width and height are set
% per figure by scripts/build_figures.py --target paper, not here: the
% paper target carries its own geometry the way the thesis and slide
% targets do, and repeating it here would silently override nothing
% while misleading the next reader.
\pgfplotsset{
  agrplot/.style={
    tick label style={font=\scriptsize},
    label style={font=\small},
    title style={font=\small\bfseries},
    legend style={font=\scriptsize, draw=none, fill=none},
    grid=major,
    grid style={gray!25, very thin},
    axis line style={gray!60},
    tick align=outside,
    tick style={gray!60},
  }}

% Ragged table cells. Justification spaces badly on a column this narrow,
% and a paper column is narrower than the thesis's.
\newcolumntype{L}[1]{>{\raggedright\arraybackslash}p{#1}}

% Room for the third pass to loosen a line rather than overrun it. The
% names in this literature are long and several are already hyphenated
% ("Think-on-Graph~\cite{...}" set 5.15pt overfull on its own), and a
% 1.5-spaced measure gives TeX fewer break points than the thesis's.
% The class ships no emergencystretch at all. 1em, not more: at 2em the
% third pass starts loosening lines it could have left alone, and the
% run trades its last overfull box for an underfull one.
\emergencystretch=1em

% A preference, not a repair. Breaking at a hyphen this literature
% already contains -- "Reasoning-" / "on-Graphs" -- reads as the start of
% a new word and costs the reader a re-parse, so it is priced high while
% ordinary hyphenation stays cheap. Do not expect it to clear an overfull
% box: swept from 0 to 4000 against the related-work paragraph, the box
% count did not move at any value. Lines that will not fit are fixed by
% rewriting the sentence, which is what happened there.
\hyphenpenalty=400
\exhyphenpenalty=4000

% Short names never break: left to itself TeX splits WebQSP after the
% "We", which a reader has to re-parse. Every entry must be letters only
% -- a digit raises "Not a letter", which is why Neo4j is absent and does
% not need to be here.
\hyphenation{WebQSP CWQ GraphRAG Freebase LangGraph McNemar Cypher AGR}

% ComplexWebQuestions is the exception and must NOT join the list above.
% At twenty letters it is longer than the slack in a 1.5-spaced
% measure, and forbidding every break in it is what set the sentence
% naming both benchmarks 46pt overfull. Hyphens here mark the breaks that
% ARE allowed, so it may split at its own word boundaries and nowhere
% else.
\hyphenation{Complex-Web-Questions}

% Consistent names for the systems, so a rename touches one line rather
% than forty. The thesis writes these out longhand in every table.
\newcommand{\noret}{no-retrieval}
\newcommand{\vecrag}{Vector-RAG}
\newcommand{\graphrag}{GraphRAG}
\newcommand{\tog}{Think-on-Graph}
\newcommand{\agr}{AGR}

% NO LINE NUMBERS, DELIBERATELY. `lineno' was loaded here and activated
% after \end{frontmatter}; both are gone, and the justification
% originally written in this comment did not survive checking. It claimed
% "Elsevier's own guidance asks for numbered lines at submission", which
% is false: Elsevier's generic guide for authors says that at initial
% submission "there are no strict formatting requirements", and does not
% mention line numbering at all (checked August 2026). The KBS guide
% specifically is still unread -- it 403s to an automated fetch by every
% route -- so a journal-level requirement cannot be ruled out; see
% journal/README.md. Do not re-add on the strength of the old
% comment, which also asserted, without checking, that Editorial Manager
% builds the reviewer PDF itself.
%
% If a submission checklist ever does ask for them, it is two lines:
% \usepackage{lineno} here, and \linenumbers after \end{frontmatter} in
% journal_0421052099.tex -- in that order and in those two places.
% Activating in the preamble instead numbers the title, author line and
% abstract heading, each on its own inconsistent count before the body
% starts.
%
% What `review' actually does is \gdef\@blstr{1.5}: 1.5 line spacing, and
% nothing else. It never provided line numbers, whatever the name suggests.

% Loaded last, which is hyperref's requirement: it redefines a great many
% internals and wants everything else already in place. `hidelinks' keeps
% the coloured boxes out of a manuscript that will be read on paper as
% often as on screen, while \Cref, \ref and \cite stay clickable for a
% reviewer who is navigating a PDF. It also supplies \url, which
% elsarticle does not.
\usepackage[hidelinks]{hyperref}

% The PDF's /Author and /Title strings are built from the raw arguments
% of \author and \title, markup and all: elsarticle accumulates them
% into \useelstitle and \useauthors (class lines 121--125 and 239--284)
% and hands both to hyperref at \begin{document} (287--290). \corref{cor1}
% is markup -- a footnote marker for the printed page -- and a metadata
% string has no footnotes, so \pdfstringdef stopped on it four times at
% \end{frontmatter}: "Token not allowed in a PDF string", for \corref
% itself, for the \cnotenum and \@corref it expands to, and for the
% \ifcase inside those.
%
% The string it produced was already right ("Md. Sakif Khan; Sadia
% Sharmin; ") because hyperref drops what it cannot use and carries on,
% so this was noise rather than damage -- but it is noise that has to be
% read and dismissed on every build, and it is what a log check counts.
%
% \pdfstringdefDisableCommands is hyperref's hook for precisely this:
% definitions that hold only while a PDF string is being built and never
% reach the page, so the asterisk still marks the corresponding author
% in the typeset frontmatter. All three of elsarticle's title-block
% reference markers are named, not just the one in use: \fnref belongs
% inside \author exactly as \corref does, and the class's own gobbling
% of \tnoteref covers the typeset title but not the copy kept for the
% metadata. Anything new inside \author or \title that expands to a mark
% rather than to text belongs in this list.
\pdfstringdefDisableCommands{%
  \renewcommand{\corref}[1]{}%
  \renewcommand{\fnref}[1]{}%
  \renewcommand{\tnoteref}[1]{}%
}

% \url breaks only at punctuation, and the repository URL in the
% data-availability statement has none in the right place: it set 0.68pt
% overfull in the 1.5-spaced measure, which survives however much
% prose surrounds it. xurl lets a URL break anywhere, which is the only
% thing that helps when the unbreakable run is "agentic-graph-reasoning".
% Load it after hyperref -- it patches hyperref's \url.
\usepackage{xurl}

% \Cref throughout, matching the thesis, which gets cleveref from
% buetcsepgthesis.sty. elsarticle does not load it either. cleveref must
% come after hyperref, and last of all, because it redefines the
% referencing commands every other package may have touched.
\usepackage{cleveref}
